% LaTeX's \lefteqn, which MathJax lacks: its argument in display style, in
% a box of no width that starts where the box is. An eqnarray row with a
% long left side starts with it, and the next rows are set from the second
% column, as in the PDF. Also in eqnarray* with \left( and \right. across
% rows, in align*, on a numbered row with a \label, in a theorem, in the
% third column, with a macro of the document inside, and in inline math.
\documentclass{article}
\usepackage{amsmath,amsthm}
\usepackage{hyperref}
\newtheorem{theorem}{Theorem}
\newcommand{\pair}[2]{\langle #1, #2 \rangle}
\begin{document}
\section{S}
\begin{eqnarray}
\lefteqn{a + b + c + e + f + g} \nonumber \\
  & = & d
\end{eqnarray}
\begin{eqnarray*}
\lefteqn{\left( \sum_{i=1}^n x_i + y \right.} \\
  & & \left. {} + z \right) = w
\end{eqnarray*}
\begin{align*}
\lefteqn{p + q + r + s} \\
  & = t
\end{align*}
\begin{eqnarray}
\lefteqn{\frac{u}{v} + \pair{u}{v}} \label{eq:long} \\
  & = & k
\end{eqnarray}
See \eqref{eq:long} and \ref{eq:long}.
\begin{theorem}
\begin{eqnarray}
h & = & \lefteqn{m + n + o + o + o} \nonumber \\
  & & {} + l
\end{eqnarray}
\end{theorem}
Inline, $\lefteqn{\sum_{j=1}^n j}xy$ overlaps.
\end{document}
